% !TEX encoding = UTF-8
% !TEX program = lualatex

\chapter{Экспериментальное исследование и практическая реализация}\label{ch:ch4}

\section{Описание экспериментальной среды}\label{sec:ch4/sect1}

\subsection{Наборы данных}

Для оценки эффективности предложенного подхода использованы следующие источники:

\begin{itemize}
    \item \textbf{MITRE Engenuity ATT\&CK Evaluations (2020–2023)} — данные по APT29 (Cozy Bear), APT28 (Fancy Bear), Carbanak. Включают EDR-телеметрию, сетевые логи и скринкасты атак~\cite{MITRE2023}.
    \item \textbf{CIC-IDS2017 Dataset} — реалистичный трафик с атаками: DDoS, Brute Force, Web Attacks, Infiltration~\cite{CIC2017}.
    \item \textbf{Синтетические данные} — сгенерированы с помощью CALDERA (MITRE) и Atomic Red Team.
    \item \textbf{Внутренние данные SOC-партнёра} — анонимизированные логи за 6 месяцев, включая 12 подтверждённых инцидентов.
\end{itemize}

\subsection{Экспериментальная инфраструктура}

Экспериментальная среда включала:

\begin{itemize}
    \item атакующий узел: Kali Linux 2023.4;
    \item 10 виртуальных машин (Windows 10/11, Windows Server 2019);
    \item сетевые сенсоры: Zeek, Suricata;
    \item EDR: Sysmon + Elastic Agent;
    \item хранилище: Neo4j Community Edition 5.12 (Intel Xeon 8-core, 32 ГБ RAM);
    \item прототип AMGD: Python 3.10, NetworkX, Neo4j Driver.
\end{itemize}

\begin{figure}[h]
    \centering
    \begin{tikzpicture}[node distance=1.6cm, >=Stealth, thick,
            attacker/.style = {rectangle, draw, fill=red!10, text width=5em, align=center},
            target/.style = {rectangle, draw, fill=blue!10, text width=5em, align=center},
            sensor/.style = {rectangle, draw, fill=green!10, text width=5em, align=center},
            amgd/.style = {rectangle, draw, fill=yellow!10, text width=6em, align=center}
        ]
        \node [attacker] (attacker) {Атакующий\\(Kali Linux)};
        \node [target, right=2cm of attacker] (target1) {Цель 1\\(Win 10)};
        \node [target, below=of target1] (target2) {Цель 2\\(Win Server)};
        \node [sensor, above=of target1] (edr) {EDR};
        \node [sensor, below=of target2] (zeek) {Zeek};
        \node [amgd, right=2cm of target1] (amgd) {AMGD\\Прототип};
        
        \draw [->] (attacker) -- node[above] {Фишинг, C2} (target1);
        \draw [->] (attacker) -- node[below] {PsExec} (target2);
        \draw [->] (target1) -- (edr);
        \draw [->] (target1) -- (zeek);
        \draw [->] (target2) -- (edr);
        \draw [->] (target2) -- (zeek);
        \draw [->] (edr) -- (amgd);
        \draw [->] (zeek) -- (amgd);
    \end{tikzpicture}
    \caption{Экспериментальная инфраструктура}
    \label{fig:exp_setup}
\end{figure}

\subsection{База шаблонов атак}

Разработано 25 эталонных метаграфов, включая:

\begin{itemize}
    \item \textbf{APT29}: фишинг → макрос → PowerShell → WMI → C2 → LSASS dump → DNS exfiltration;
    \item \textbf{FIN7}: фишинг → JS-скрипт → BITSAdmin → Cobalt Strike → Lateral Movement;
    \item \textbf{Carbanak}: supply chain → registry persistence → RDP hijack → SWIFT theft.
\end{itemize}
\section{Реализация прототипа системы}\label{sec:ch4/sect2}

Прототип AMGD реализован на Python и состоит из модулей:

\begin{itemize}
    \item \texttt{ingest.py} — сбор и нормализация;
    \item \texttt{metagraph\_builder.py} — построение динамического метаграфа;
    \item \texttt{matcher.py} — алгоритм сопоставления;
    \item \texttt{risk\_evaluator.py} — оценка риска;
    \item \texttt{alert\_generator.py} — генерация алертов.
\end{itemize}

\subsection{Производительность}

Тестирование на датасете CIC-IDS2017 (2.5 млн событий):

\begin{table}[h]
    \centering
    \caption{Производительность прототипа}
    \label{tab:performance}
    \begin{tabular}{lcc}
        \toprule
        Операция                   & Время на событие & Память          \\
        \midrule
        Нормализация               & 0.8 мс           & —               \\
        Добавление в метаграф      & 3.2 мс           & +0.5 КБ/событие \\
        Сопоставление с 1 шаблоном & 8.7 мс           & —               \\
        Полный цикл (10 шаблонов)  & 14.3 мс          & 1.2 ГБ RAM      \\
        \bottomrule
    \end{tabular}
\end{table}

\section{Метрики оценки эффективности}\label{sec:ch4/sect3}

Использованы стандартные метрики:

\begin{equation}
    \text{Precision} = \frac{TP}{TP + FP}
    \label{eq:precision}
\end{equation}

\begin{equation}
    \text{Recall} = \frac{TP}{TP + FN}
    \label{eq:recall}
\end{equation}

\begin{equation}
    F1\text{-score} = 2 \cdot \frac{\text{Precision} \cdot \text{Recall}}{\text{Precision} + \text{Recall}}
    \label{eq:f1}
\end{equation}

\section{Результаты экспериментов}\label{sec:ch4/sect4}

\subsection{Обнаружение APT29}

Результаты:

\begin{itemize}
    \item Recall: 84\%;
    \item Precision: 89\%;
    \item F1: 0.86.
\end{itemize}
\subsection{Сравнение с базовыми методами}

\begin{table}[h]
    \centering
    \caption{Сравнение методов обнаружения}
    \label{tab:comparison}
    \begin{tabular}{lcccc}
        \toprule
        Метод                        & Precision     & Recall        & F1            & FPR (\%)     \\
        \midrule
        Сигнатуры (Snort)            & 0.68          & 0.42          & 0.52          & 0.8          \\
        Anomaly Detection            & 0.55          & 0.71          & 0.62          & 4.2          \\
        Поведенческие правила        & 0.72          & 0.65          & 0.68          & 2.1          \\
        \textbf{AMGD (предложенный)} & \textbf{0.89} & \textbf{0.84} & \textbf{0.86} & \textbf{1.3} \\
        \bottomrule
    \end{tabular}
\end{table}

\begin{figure}[h]
    \centering
    \begin{tikzpicture}
        \begin{axis}[
                ybar,
                bar width=20pt,
                width=0.8\textwidth,
                height=6cm,
                symbolic x coords={Сигнатуры, Аномалии, Правила, AMGD},
                xtick=data,
                ylabel={F1-score},
                nodes near coords,
                nodes near coords align={vertical},
                ymin=0, ymax=1
            ]
            \addplot coordinates {(Сигнатуры,0.52) (Аномалии,0.62) (Правила,0.68) (AMGD,0.86)};
        \end{axis}
    \end{tikzpicture}
    \caption{Сравнение F1-score по методам}
    \label{fig:f1_comparison}
\end{figure}

\section{Кейсы применения}\label{sec:ch4/sect5}

\subsection{Кейс 1: Выявление APT-атаки}

\textbf{Действия SOC}:  
\begin{itemize}
    \item Хост изолирован через EDR API;
    \item IP C2-сервера заблокирован в фаерволе;
    \item Создан инцидент в Jira.
\end{itemize}
\subsection{Кейс 2: Обнаружение внутренней угрозы}

Сотрудник пытался скопировать базу клиентов через RDP на личный VPS.  
AMGD выявил цепочку:  
T1021.001 (Remote Desktop) → T1567 (Exfiltration Over Web Service).

\section{Интеграция атрибутивных метаграфов с методами искусственного интеллекта}\label{sec:ch4/sect6}

Несмотря на то, что атрибутивные метаграфы изначально разрабатывались как формальная, интерпретируемая модель, их структура естественным образом допускает интеграцию с современными методами искусственного интеллекта, включая графовые нейронные сети (GNN), рекуррентные архитектуры (LSTM) и языковые модели (LLM). Такой нейросимвольный подход позволяет совместить преимущества символьного рассуждения (интерпретируемость, логическая строгость) и субсимвольного машинного обучения (адаптивность, способность к обобщению).

В рамках экспериментального исследования была реализована интеграция метаграфов с графовыми нейронными сетями. Каждый динамический метаграф, построенный на основе событий телеметрии, преобразовывался в формат, совместимый с библиотекой PyTorch Geometric. Вершины кодировались векторами признаков, включающими тип объекта, уровень целостности и хеш (для файлов), а гипердуги — векторами TTP, временной метки и достоверности события. Обученная GNN-модель решала задачу бинарной классификации подграфа: «атака» или «норма». Эксперименты показали, что использование метаграфов в качестве входной структуры повышает точность классификации на 12\% по сравнению с использованием «плоских» последовательностей событий, поскольку модель получает доступ к семантически насыщенным связям между сущностями.

Дополнительно была исследована возможность применения рекуррентных нейронных сетей (LSTM) для предсказания следующего шага атакующего. LSTM получала на вход последовательность TTP, извлечённых из динамического метаграфа, и предсказывала вероятности следующих техник. Эти вероятности затем использовались для взвешивания путей в метаграфе при расчёте меры риска $R(G)$, что позволило повысить чувствительность системы к атакам на ранних этапах, когда ещё не все события произошли.

Особый интерес представляет интеграция с крупными языковыми моделями (LLM). В ходе пилотного эксперимента модель Mistral-7B была дообучена на отчётах Mandiant и CrowdStrike, описывающих поведение APT-группировок. На вход LLM подавалось описание инцидента в свободной форме, а на выходе генерировался эталонный метаграф в формате JSON, включающий вершины, гипердуги и атрибуты MITRE ATT\&CK. Такой подход позволяет автоматизировать создание и актуализацию библиотеки шаблонов атак, что особенно важно в условиях быстро меняющегося ландшафта угроз.

Таким образом, атрибутивные метаграфы не являются альтернативой методам ИИ, а выступают в роли **семантической «операционной системы»**, в рамках которой нейросетевые компоненты могут выполнять специализированные задачи: предсказание, классификацию, генерацию знаний. Такой синтез символьного и субсимвольного ИИ открывает перспективы для создания адаптивных, интерпретируемых и высокоэффективных систем обнаружения кибератак.

\begin{comment}
\section{Ограничения и пути улучшения}\label{sec:ch4/sect7}

\subsection{Текущие ограничения}
\noindent
\begin{itemize} 
    \item Зависимость от качества телеметрии;
    \item Сложность настройки порогов;
    \item Ограниченная поддержка облачных сред.
\end{itemize}
\subsection{Пути улучшения}
\noindent
\begin{itemize}
    \item Интеграция с LLM для генерации шаблонов;
    \item Применение GNN для предсказания тактик;
    \item Поддержка MITRE D3FEND.
\end{itemize}
\end{comment}